% Review 2010-08-13: NASA-Bild kontrovers
% \begin{savequote}
% \includegraphics[width=5cm]{./images/nasa-pd/157684main_image_feature_656_ys_full-00800.jpg}
% 
% Bild: NASA
% \end{savequote}
\def\ChapterCode{GYN}
\chapter
[\CO{[GYN]} \emph{Spezielle Patientengruppe:} Frauen und Schwangere]
{\emph{Spezielle Patientengruppe:} Frauen und Schwangere}

\ChapterIntro
{%
~~~
}
{%
\begin{description}
  \item[Maintainer:] Sebastian Gabriel
  \item[Autoren:] Diverse
  \item[Reviewer:] Standard-Reviewprozess
  \item[Version:] {\VarDocVersion} ({\VarDocVersionDescription})
  \item[SHA1:] {\DocGitCommit}
\end{description}

{\TxTeilkapitelAASS}
}


\clearpage % TEMP temp clearpage
\section{Schwangerschaft}
\HeadingSub{\Lang{Term.} \I{Gestation}, \I{Gravidität}, \Lang{Abkz.} \I{Grav.}}

\subsection{Allgemeines}

\TextSlide{Exkurs: Weiblicher Zyklus}{%
Der Begriff Zyklus oder Periode beschreibt in diesem
Zusammenhang den Vorgang des rund alle \EE{28 Tage}
stattfindenden Reifens mindestens einer Eizelle, dem
Eisprung und den damit im Zusammenhang stehenden
Veränderungen der Gebärmutterschleimhaut. Der Zyklus ist
unter
\FullPrettyRef{topic:weiblicher-zyklus}, beschrieben.
}
{\FullPrettyRef{topic:weiblicher-zyklus}}
{}{}{}









\TextSlide{Schwangerschaftsdauer}{%
Die normale Schwangerschaft dauert \EE{40\,$\pm$\,2
Wochen}. Das sind rund \EE{9 Kalendermonate} \Z{oder 10
Lunarmonate\footnote{\Z{\Lang{lat.} luna: Mond.}}}. Der
Fortschritt der Schwangerschaft wird in
\T{Schwangerschaftswochen}, Abk. \T{SSW}, angegeben. Häufig
wird auch die Unterteilung in drei
\T[Trimenon]{Trimester}\ZFootnote{\Lang{lat.} tri: drei; \Lang{lat.}
mensis: Monat; Trimenon, Trimester: dreimonatig} zu je 3
Kalendermonaten verwendet.
Von
einer \T{Frühgeburt} spricht man wenn die Geburt
\E{bis zur vollendeten 37. SSW}
erfolgt. Von einer \E{Übertragung} spricht man ab der
\E{42.~SSW}.
}{%
\begin{itemize}
  \item 40 $\pm$ 2 Schwangerschaftswochen (SSW)
  \item 3 Trimester
  \item Frühgeburt: bis 37. SSW
  \item Übertragung: > 42. SSW
\end{itemize}
}{}{}{}


\TextSlide{Schwangerschaftsverlauf}
{%
Nach dem Eisprung ist die Eizelle \VonBis{8}{10} Stunden
befruchtungsfähig. Sie benötigt für die Wanderung Eierstock
\Sign{→} Eileiter \Sign{→} Uterus (Gebärmutter)
\VonBis{4}{6} Tage. Wurde die Eizelle befruchtet, kommt es
zur \EE{Einnistung in der Gebärmutterschleimhaut}. Dort
bildet sich der \T{Mutterkuchen} (\T{Plazenta}), welcher 
das Blut der
Frucht\footnote{Mit dem Begriff \Speech{\enquote*{Frucht}}
ist das Kind gemeint. Die  Frucht hat je nach
Schwangerschaftsfortschritt verschiedene Namen: Man spricht
vom \I{Embryo} bis zur 12. SSW, danach vom \I{Fetus}.} mit
Nährstoffen und Sauerstoff versorgt.
Weiters bildet sich um die Frucht eine \T{Fruchtblase},
welche mit \T{Fruchtwasser} gefüllt ist und sie so ideal vor
äußeren Erschütterungen schützt.

Zwischen der 18. und der 20. SSW kann die Mutter die ersten
Kindesbewegungen spüren. \EE{Ab der 24.
SSW gilt das Kind als überlebensfähig} und ggfs. bei einer
Frühgeburt als reanimationspflichtig.
}
{%
\begin{itemize}
  \item Eierstock \Sign{→} Eileiter \Sign{→} Gebärmutter
  \item Befruchtung
  \item Einnistung in Gebärmutterschleimhaut
  \item Mutterkuchen (Plazenta)
  \item Fruchtblase, Fruchtwasser
  \item Ab 24. SSW überlebensfähig
\end{itemize}

}
{}
{}
{}

\Text
{Begriffe und gängige Abkürzungen}
{
\begin{description}
  \item[Gravida] \Lang{Syn.} Grav., Gravididät.
  Schwangerschaft, bzw. Anzahl der bisherigen
  Schwangerschaften inklusive der aktuellen. Z.\,B. Grav
  III: 3. Schwangerschaft
  \item[\Z{Para}] \Z{Anzahl der vorangegenagenen Geburten.
  Sonderfall \I{Nullipara}: Noch keine Geburt}
  \item[Partus] Wievielte Geburt
  \item[Abort] Fruchtabgang
  \item[\Z{Interruptio}] \Z{Künstliche Beendigung der
  Schwangerschaft}
  \item[LNR] \Lang{Abkz.} Letzte normale Regelblutung
  \item[SSW] Schwangerschaftswoche, Angabe in vollendeten
  Schwangerschaftswochen mit dem Zusatz der Tage der
  angebrochenen Schangerschaftswoche, z.\,B. SSW 38+4
\end{description}

\subparagraph{Beispiel} \Speech{\enquote{Grav. IV SSW 38+4
Partus II Abort II}}: Patientin in der 4. Schwangerschaft am
4.
Tag der 39. Schwangerschaftswoche, zweitgebärend (Par\underline{a} I bzw. Part\underline{us} II) , zwei Fruchtabgänge.

}
{}
{}
{}
{}

\begin{PictureSeriesWide}{}{Bilderserie: Schwangerschaft}
\PictureSeriesRowWide{3}
{./images/wikipedia-pd/Nabelschnur.png}
{Die Lage des Fötus in der Gebärmutter.
\SourcePicture{Gray's Anatomy, 20th ed., PD, Copyright
expired.}}
{empty}
{}
{empty}
{}
{empty}
{}
\end{PictureSeriesWide}




\subsection{Notfälle -- Frühschwangerschaft}
\label{topic:notfall-fruehschwangerschaft}

\subsubsection{Fehlgeburt}
%
\index{Fehlgeburt}%
\index{Abort}%
\label{topic:fehlgeburt}%
\label{topic:abort}%

\DefinitionInPara{Ein \T{Abort} ist die Beendigung einer Schwangerschaft 
bzw. Fruchtabgang bei einem
Kindesgewicht von unter 500\Gramm* und fehlenden Lebenszeichen.
\cite{HaagGyn2007,ForMed2}}


\TextSlide{\TxBeschreibung}{%
Aufgrund verschiedener Ursachen kann es zum Absterben der
Frucht mit einer nachfolgenden Fehlgeburt (Ausstoßung)
kommen. Wenn dies vor dem Erreichen von 500\Gramm*
geschieht, spricht man von einer \T{Fehlgeburt}.
}{
\begin{itemize}
  \item Frühe Fehlgeburt\,/\,Fruchtabgang
  \item Oft unbemerkt
\end{itemize}
}{}{}{}


\TextSlide{\TxSymptome}{%
Die Symptome sind vor allem davon abhängig, wie weit die
Schwangerschaft fortgeschritten ist bzw. die Frucht
entwickelt ist. Ein sehr früher Fruchtabgang wird oft nicht
einmal bemerkt bzw. mit Regelunregelmäßigkeiten verwechselt,
oft weiß die Frau nicht einmal, dass sie schwanger gewesen
ist.  Ein Abort macht sich dann vor allem durch eine
\EE{vaginale Blutung} bemerkbar, die mitunter als
\Speech{\enquote{verstärkte Regelblutung}} fehlgedeutet wird.
Weiters kann es zum \EE{Abgang von Gewebeteilen} und
\EE{Unterbauchschmerzen} kommen.
Eine Fehlgeburt kann jedoch auch zu einem späteren Zeitpunkt
während der Schwangerschaft vor Erreichen der
500\Gramm*-Grenze des Kindes vorkommen.\ZFootnote{Nach dem
Erreichen der 500\Gramm*-Marke spricht man von einer
\E{Totgeburt}; für diese gelten z.\,T. andere rechtliche
Rahmenbedingungen (Kind kann einen Vornamen erhalten,
Beurkundung im Sterbebuch).
\cite{ForMed2}.}
}{
\begin{itemize}
	\item Vaginale Blutung
	\item Abgang von Gewebsteilen
	\item Unterbauchschmerzen
\end{itemize}
}{}{}{}

\TextSlide{Anmerkung}{% 
Sollten Zweifel bestehen, ob es sich um einen Abort handelt,
sollte die Verdachtsdiagnose \Speech{\enquote{V.\,a. Abort}}
\EE{vermieden} werden, um die Patientin nicht --
möglicherweise unnötig -- zu beunruhigen. Stattdessen sollte
man eine korrekte Symptombeschreibeung als Verdachtsdiagnose
verwenden, z.\,B. \Speech{\enquote{Vaginale Blutung in 5. SSW}}.

Sollten Zweifel hinsichtlich der
Lebensfähigkeit des Kindes bestehen, muss mit einer
Reanimation begonnen werden.

\Cave{Im Zweifel: Reanimation!}
}{% 
Im Zweifel reanimieren! 

Diagnose \enquote{Abort} vermeiden \Sign{→} Symptombeschreibung}{}{}{}

% M %%%%%%%%%%%%%%%%%%%%%%%%%%%%%%%%%%%%%%%%%%%%%%%%%%% M %
% Abort
\ProcedurePrintDefault{GO06XX1C}



\subsubsection{Eileiterschwangerschaft}
\label{topic:extrauterine-graviditaet}

\TextSlide{\TxBeschreibung}{%
\index{Eileiterschwangerschaft}%
\DefinitionInPara{Eine \T{Eileiterschwangerschaft} ist eine
Schwangerschaft, bei der sich die Frucht im Eileiter
einninstet. {\Lang{Term.} {Tubaria}, \Lang{Syn.}
extrauterine Schwangerschaft.}}
Das Problem dabei: Mit zunehmender Größe der Frucht wird der
Eileiter gedehnt, dabei kommt es zu Schmerzen. Schließlich reißt der Eileiter
und die Frau erleidet massive innere Blutungen. Es zeigen
sich dann Anzeichen eines akuten Abdomens und eines
hypovolämischen Schocks.
}{
\begin{itemize}
    \item Einnistung der Frucht in Eileiter (statt Gebärmutter)
    \item Schmerz durch Dehnung des Eileiters
    \item $\rightarrow$ Zerreißung des Eileiters
    \begin{itemize}
        \item Blutung
        \item Akutes Abdomen
        \item Schock
    \end{itemize}
\end{itemize}
}{}{}{}


\Fazit{Bei Frauen mit akutem Abdomen muss immer eine mögliche
Schwangerschaft abgeklärt werden!}

% M %%%%%%%%%%%%%%%%%%%%%%%%%%%%%%%%%%%%%%%%%%%%%%%%%%% M %
% Eileiterschwangerschaft, Verdacht
\ProcedurePrintDefault{GO00091C}


% M %%%%%%%%%%%%%%%%%%%%%%%%%%%%%%%%%%%%%%%%%%%%%%%%%%% M %
% Eileiterschwangerschaft, rupturiert
\ProcedurePrintDefault{GO00092C}

% \clearpage % TEMP temp clearpage
\subsection{Notfälle -- Spätschwangerschaft}
\label{topic:notfall-spaetschwangerschaft}

\subsubsection{Vorzeitige Plazentalösung}

\TextSlide{\TxBeschreibung}{%
\DefinitionInPara{Eine
\T[Plazentaöslung!vorzeitige]{vorzeitige Plazentalösung} ist
eine Ablösung des Mutterkuchens (Plazenta) vor Beendigung der Geburt}
Aufgrund einer Einblutung zwischen Plazenta und Gebärmutter
kann es zu einer vorzeitigen Ablösung der Plazenta mit dann
ungehinderter \EE{Blutung} kommen. Es besteht Schock- bzw.
\EE{Lebensgefahr} für Mutter und Kind durch Blutung der Mutter und
Mangelversorgung des Kindes.
Die Diagnose ist präklinisch nicht sicher zu
stellen. Im Vordergrund steht daher die Behandlung der
Symptome und ein zügiger Transport.
}
{
\begin{itemize}
  \item Blutung
  \item Schock, Lebensgefahr
\end{itemize}
}
{}{}{}


\TextSlide{\TxSymptome}{%
Die Symptome sind unspezifisch: Im Vordergrund steht der
\EE{Schock} infolge der Blutung. Zusätzlich zu den
Schocksymptomen kann die Frau Symptome eines \EE{akuten
Abdomens} zeigen.

\bigskip

\bigskip
}{%
\begin{itemize}
  \item Unspezifisch:
  \begin{itemize}
    \item Schock
    \item Akutes Abdomen
  \end{itemize}
\end{itemize}
}{}{}{}



% M %%%%%%%%%%%%%%%%%%%%%%%%%%%%%%%%%%%%%%%%%%%%%%%%%%% M %
% Vorzeitige Plazentalösung
\ProcedurePrintDefault{GO45090C}

\subsubsection{Vena-cava-Kompressionssyndrom}
\label{topic:vena-cava-syndrom}
\index{Vena-Cava-Kompressionsyndrom}

\Text{\TxBeschreibung}{%
Bei hochschwangeren Frauen kann es vorkommen, dass das
\EE{Kind die untere Hohlvene\index{Hohlvene!untere!Vena-Cava-Kompressionssyndrom} 
der Mutter abdrückt}. Dadurch
wird der Blutstrom zum Herzen vermindert und es kann zu
einem \EE{plötzlichen Kollaps der Mutter} und einer
Unterversorgung des Kindes kommen. Durch eine
Linksseitenlagerung wird der Druck, den das Kind ausübt,
vermindert, da die Hohlvene eher rechts der Körpermitte
verläuft.


}{

}
{
\begin{itemize}
  \item Kind drückt untere Hohlvene der Mutter ab
  \item Kollaps
\end{itemize}
}
{}
{}

\Fazit{\EE{Jede hochschwangere Patientin} soll in (leichter)
Linksseitenlage und nicht in Rückenlage transportiert werden. }

% M %%%%%%%%%%%%%%%%%%%%%%%%%%%%%%%%%%%%%%%%%%%%%%%%%%% M %
% Vena-cava-Syndrom
\ProcedurePrintDefault{GO26051C}


\subsubsection{Präeklampsie, EPH-Gestose, HELLP-Syndrom, Eklampsie}
\label{topic:eph-gestose}
\label{topic:praeeklampsie}
\label{topic:eklampsie}
\label{topic:hellp-syndrom}
\index{EPH-Gestose}
\index{Präeklampsie}
\index{Eklampsie}
\index{HELLP-Syndrom}
\index{Krampfanfall!eklamptischer}

% FEATURE INHALT: HELLP,Syndrom, Leberspannungsschmerz

\TextSlide{\TxBeschreibung}{%
Die \T{EPH-Gestose}, auch \T{Präeklampsie} genannt,
ist eine gefürchtete Komplikation der
Spätschwangerschaft, bei der es zu einer
Durchblutungsstörung der Gebärmutter, und in weiterer Folge
zu diversen Komplikationen bei der Mutter kommt. Sie wird
in der Regel im Rahmen der Mutter-Kind-Pass-Untersuchungen
erkannt. Die
Maximalvariante der Präeklampsie ist das \T{HELLP-Syndrom}.
Unbehandelt kann die Präeklampsie und das HELLP-Syndrom zu
einem \E{eklamptischen Krampfanfall} (\T{Eklampsie}, ähnlich
einem \enquote*{normalen} zerebralen Krampfanfall) und in
weiterer Folge zum Tod von Mutter und Kind führen.
Die Ursachen sind nicht vollständig geklärt.
}
{
\begin{itemize}
  \item Schwangerschaftserkrankung
  \item Kann zur Eklampsie (Krampfanfall) führen
\end{itemize}
}
{}
{}
{}

\Fazit{Die \emph{Prä}eklampsie (EPH-Gestose) kann zu einem
lebensgefährlichen Krampfanfall (Eklampsie) führen.}

\TextSlide{\TxSymptome}{
Die Präeklampsie ist charakterisiert durch:
\begin{enumerate}
    \item \EE{E}dema (Ödeme)
    \item \EE{P}roteinuria (Proteinurie = Eiweiß im Harn, Harn schäumt)
    \item \EE{H}ypertonie (RR  ${\uparrow}$) 
    
    Ein Blutdruck
    {\textgreater}\,\RR{140}{90} bei einer Schwangeren ist zu  hoch!
\end{enumerate}
Beim HELLP-Syndrom kommt es zu einem \EE{rechten
Oberbauchschmerz} (Leberdehnungsschmerz).
Bei einer drohenden Eklampsie klagt die Patientin über 
\EE{Kopfschmerzen}, Sehstörungen sowie Übelkeit und
Erbrechen. \cite{Haag2006}

}{
\begin{itemize}
    \item \EE{E}dema (Ödeme)
    \item \EE{P}roteinuria (Proteinurie = Eiweiß im Harn)
    \item \EE{H}ypertonie (RR  ${\uparrow}$)
    \item HELLP: re. Oberbauchschmerz
    \item Drohende Eklampsie: \begin{compactitem}
      \item Kopfschmerzen
      \item Sehstörungen
      \item Übelkeit, Erbrechen
    \end{compactitem}
\end{itemize}
}{}{}{}


\SlideCompensationNoParskip{1}





% M %%%%%%%%%%%%%%%%%%%%%%%%%%%%%%%%%%%%%%%%%%%%%%%%%%% M %
% Präeklampsie
\ProcedurePrintDefault{GO14090C}


% M %%%%%%%%%%%%%%%%%%%%%%%%%%%%%%%%%%%%%%%%%%%%%%%%%%% M %
% Eklampsie, eklamptischer Krampfanfall
\ProcedurePrintDefault{GO15091C}


\subsubsection{Vorzeitiger Fruchtwasserabgang}

\Text{\TxBeschreibung}{%
Der Fötus schwimmt während  Schwangerschaft in steriler
Flüssigkeit (Fruchtwasser). Bei vorzeitiger Eröffnung der
Fruchtblase ergeben sich für ihn die Gefahr
eines \EE{Nabelschnurvorfalls} oder einer \EE{Infektion}
(Keimbesiedlung der Fruchthöhle). Die Patientin ist an
einer geburtshilflichen Abteilung vorzustellen.
}{}{}{}{}

% M %%%%%%%%%%%%%%%%%%%%%%%%%%%%%%%%%%%%%%%%%%%%%%%%%%% M %
% Vorzeitiger Blasensprung
\ProcedurePrintDefault{GO42090C}


\clearpage
\section{Geburtshilfe}

\subsection{Die Geburt}
\label{topic:geburt}

Als \EE{Beginn der Geburt} zählt:
\begin{itemize}
  \item Regelmäßige Wehen in \VonBis{10}{15} Minuten
  Abstand auftreten wenn und dieser Zustand länger als eine
  Stunde dauert 
  \item Sprung der Fruchtblase und Abgang von
  Fruchtwasser 
\end{itemize}
Die Geburt verläuft in 3 Phasen: 
\begin{itemize}
  \item \T{Eröffnungsphase}
  \item \T{Austreibungsphase} 
  \item \T{Nachgeburtsphase}
\end{itemize} 


\TextSlide{Eröffnungsphase}
{%
Die Eröffnungsphase beginnt mit den ersten regelmäßigen
Wehen und endet mit der vollständigen Eröffnung des
Muttermundes auf 10\CM*.
\cite{StriebelAnaesthesie2,BuehlingGynaekologie1,HaagGyn2007}
Durch Zusammenziehen der Gebärmuttermuskulatur
(Wehen) wird das Kind langsam in den Geburtskanal nach unten
gepresst. Dadurch verkürzt sich der Gebärmutterhals
(Zervix), und der Muttermund wird eröffnet. Der
Wehenabstand beträgt ca. \VonBis{2}{10}\Min*.
Am Ende der Eröffnungsphase
kommt es zum Platzen der Fruchtblase (\T{Blasensprung}),
dabei geht Fruchtwasser und etwas blutiges Sekret ab.

\Z{Bei Erstgebärenden dauert die Eröffnungsperiode ca.
\VonBis{6}{9}\Hour*, bei Mehrgebärenden ca.
\VonBis{3}{7}\Hour* \cite{BuehlingGynaekologie1}. Die
tatsächliche Dauer kann \E{sehr} unterschiedlich sein.}

\Cave{%
Ab dem Blasensprung darf die Gebärende nicht mehr aufstehen,
sondern muss liegend transportiert werden. Es besteht die
Gefahr eines Nabelschnurvorfalls in den Geburtskanal!}
}{%
\begin{itemize}
  \item Von Beginn regelmäßger Wehen bis vollst. Eröffnung
  des Muttermundes
  \item Wehen im Abstand von \VonBis{2}{10} Minuten
  \item Gebärmutterhals verkürzt sich, Muttermund öffnet sich
  \item Blasensprung mit Fruchtwasserabgang
  \item Ab Blasensprung: \EE{Gefahr eines
  Nabelschnurvorfalles}
  \begin{itemize}
    \item \Ca{Liegender Transport!}
  \end{itemize}
\end{itemize}
}{}{}{}

\TextSlide{Austreibungsphase}
{%
Die Austreibungsphase beginnt mit der
vollständigen \E{Eröffnung} des Muttermundes \Z{(10\CM*)}
und endet mit der \E{Geburt} des Kindes.
\cite{BuehlingGynaekologie1,HaagGyn2007}
Der Wehenabstand beträgt nun \EE{\VonBis{2}{3}}
Minuten. Die Wehen werden im Verlauf kraftvoller und die
Frau verspürt mit Fortschreiten der Geburt einen
\EE{Pressdrang}(\EE{Presswehen}). Der Kopf des Kindes wird
schließlich nun zwischen den Schamlippen  sichtbar.
Wenn der Kopf geboren ist, folgt der restliche Körper in der
Regel schnell nach. Vorsicht: Das Kind ist durch die
Fruchtschmiere rutschig!

\subparagraph{Der Damm}
\label{topic:dammschutz}
Als \T{Damm} wird das Stück Haut zwischen Scheide und Anus
bezeichnet. Mittels eines
Griffes namens \T{Dammschutz} kann
theoretisch die Wahrscheinlichkeit eines Dammrisses
minimiert werden. Dabei wird die Hand unterhalb der
Scheide angelegt und man versucht das Gewebe zu
stabilisieren. In der Praxis ist dieser Schutz, insbesonders
für darin nicht routiniertes Personal, schwer zu erreichen
\cite{AsboeLehrmeinung2011-20}.

Dennoch ist der Dammschutz nicht unwichtig: Durch das
Pressen wird nicht nur das Kind geboren, sondern in der
Regel auch  \EE{Stuhl} abgesetzt. Durch den Dammschutz, in
Verbindung mit einer Unterlage, kann das Kind \E{vor dem
Stuhl der Mutter geschützt} werden (\T{Stuhlschutz}).

}{%
\begin{itemize}
  \item Von vollständiger Eröffnung des Muttermundes bis
  Geburt des Kindes
  \item Wehen im Abstand von \VonBis{2}{3}\Min*
  \item Presswehen
  \item Dammschutz = Stuhlschutz
  \item Geburt des Kindes
\end{itemize}
}
{}
{}
{---}






\TextSlide{Nachgeburtsphase}
{%
Nach \VonBis{10}{20} Minuten kommt es neuerlich zu Wehen,
welche den Presswehen ähnlich sind. Es kommt zur Geburt der
\E{Plazenta} (\E{Mutterkuchen}). Die Gefahr hierbei ist,
dass es zu starkem Blutverlust bei der Mutter kommen kann.
\EE{Die Plazenta ist für weitere
Untersuchungen\footnote{Kontrolle der Vollständigkeit,
Prüfen auf Insuffizienzzeichen, \ldots} ins
Krankenhaus mitzunehmen!} 
}{
\begin{itemize}  
    \item Nach \VonBis{10}{20} Min. neuerlich Wehen. 
    \item Geburt der Plazenta.
\end{itemize}
}{}{}{}



\Text{\TxLink}{%
\begin{itemize}
  \item Birth of Baby (Vaginal Childbirth):
\url{http://www.youtube.com/watch?v=Xath6kOf0NE}
\end{itemize}

}{}{}{}{}


\begin{PictureSeriesWide}{}{Nabelschnur}
\PictureSeriesRow{3}
{./images/wikipedia-pd/Newborn_umbilical_suction-AASS-0030mm.jpg}{Durchtrennung der Nabelschnur, die patientenferne Klemme ist durch die Hand verdeckt. \SourcePicture{WMC: \enquote{Jeremykemp} / PD}}
{./images/wikipedia-pd/Navelklem.jpg}{Abgeklemmte Nabelschnur \SourcePicture{WMC: \enquote{Harmid} / PD}}
{./images/wikipedia-pd/Umbilicalstump.jpg}{Nabelschnurstumpf nach sieben Tagen \SourcePicture{WMC: \enquote{Fatrabbit} / PD}}
{}{}
\end{PictureSeriesWide}



\subsubsection{Die bevorstehende Geburt}


\TextSlide{Anamnese -- Spezielle Fragen}
{%
\begin{itemize}
  \item Wievielte Schwangerschaftswoche
  (SSW)? Geplanter Geburtstermin?
  \item Wievielte Schwangerschaft,
  wievielte Geburt? 
  \item Wehenabstand
  \item Blasensprung? 
  \item Komplikationen bekannt?
  \item \T{Mutter-Kind-Pass} verlangen. Darin stehen mit
  unter wichtige Informationen über die Schwangerschaft bzw. die  Geburt.
  \item In welchem Krankenhaus ist die Geburt
  angemeldet?
\end{itemize}
}{
\begin{itemize}
  \item SSW, Termin?
  \item WievielteSchwangerschaft bzw. Geburt?
  \item Wehenabstand?
  \item Blasensprung?
  \item Komplikationen bekannt?
  \item Mutter-Kind-Pass
  \item Voranmeldung zur Geburt?
\end{itemize}
}{}{}{}
  
  
\Text{Betreuung der Gebärenden und Transport}
{%
Anhand des Wehenabstandes und dem Vorliegen oder
Nichtvorliegen eines Pressdrangs oder Blasensprungs kann
man den Fortschritt der Geburt einschätzen. 
Bei Mehrgebärenden ist der Geburtsverlauf oft rascher. 
Befindet
sich die Mutter noch in der Eröffnungsphase wird ein 
schnellstmöglicher, schonender Transport ins Krankenhaus 
angestrebt. Grundsätzlich soll die Gebärende in,
zumindest leichter, \EE{Seitenlage} transportiert werden
um ein \E{Vena-cava-Kompressionssyndroms}
(\FullPrettyRef{topic:vena-cava-syndrom}) zu vermeiden.
Idealerweise wird nach lins (Linksseitenlage) gelagert, doch
ist das nicht zwingend. Ab dem erfolgtem \EE{Blasensprung} darf die
werdende Mutter jedenfalls nur mehr \EE{liegend}
transportiert werden, um einen Nabelschnurvorfall vorzubeugen.
Sollte die Mutter stehend
angetroffen werden, muss dies dokumentiert werden.
}
{}



\subsubsection{Der Geburtsvorgang und Versorgung des Neugeborenen}


Ab Beginn der Presswehen bzw. wenn der Kindskopf oder -teile
(Nabelschnur, Hand, \ldots) zwischen den Schamlippen
sichtbar sind, ist der weitere Transport zu
unterlassen und das Fachpersonal muss bei der Geburt
Hilfestellung leisten.


% M %%%%%%%%%%%%%%%%%%%%%%%%%%%%%%%%%%%%%%%%%%%%%%%%%%% M %
% Spontangeburt
\ProcedurePrintDefault{GO80XX1C}


% M %%%%%%%%%%%%%%%%%%%%%%%%%%%%%%%%%%%%%%%%%%%%%%%%%%% M %
% Versorgung des Neugeborenen
\ProcedurePrintDefault{GY51310C}

\TextSlide
{Apgar-Score}
{Der \I{Apgar-Score} ist nach Virginia Apgar
benannt und ist ein Punkteschema zur Beurteilung von
Neugeborenenen. Die Bestimmung erfolgt 1, 5 und 10 Minuten
nach der Geburt. Die optimale Punktzahl für Neugeborene sind
\Range{9}{10} Punkte. \Z{Bei Wertungen zwischen \Range{5}{8}
gilt das Neugeborene als gefährdet, unter 5 als akut
lebensgefährdet. Für die weitere Behandlung ist jedoch der
klinische Zustand entscheidend, nicht der Apgar-Score.}
\Cave{Die Versorgung des Neugeborenen hat absoluten Vorrang 
vor der Erhebung des Apgar-Score!}
}
{
\begin{itemize}
  \item Punkteschema zur Beurteilung von 
  Neugeborenenen
  \item Max. \Range{9}{10} Punkte
\end{itemize}
}
{}
{}
{}

\Layout{57}
{}
{}
{}
{\Z{\begin{tabularx}{\linewidth}{XXXX}
\addlinespace\toprule\addlinespace
\noindent\TabHeading{Kriterium} & \noindent\TabHeading{0 Punkte} & \noindent\TabHeading{1 Punkt} & \noindent\TabHeading{2 Punkte}
\\\addlinespace\midrule\addlinespace
\noindent\TabSub{Herzfrequenz} & unter 60\PerMin* & \Range{60}{99}\PerMin* & über 100\PerMin*  
\\\addlinespace
\noindent\TabSub{Atemanstrengung} & keine & unregelmäßig, flach &  regelmäßig, Kind schreit
\\\addlinespace
\noindent\TabSub{Reflexe} & keine & Grimassieren &  kräftiges Schreien
\\\addlinespace
\noindent\TabSub{Muskeltonus} & schlaff & leichte Beugung der Extremitäten & aktive Bewegung der Extremitäten 
\\\addlinespace
\noindent\TabSub{(Haut-) Farbe} & blau, blass & Stamm rosig, Extremitäten blau & gesamter Körper rosig 
\\\addlinespace\bottomrule\addlinespace
\end{tabularx}
}}
{Apgar-Score}
{}






Nach dem Ende der Geburt müssen die \EE{Geburtszeit} und der
\EE{Geburtsort} (Straße mit Hausnummer; Kreuzung; Autobahn
mit km-Angabe, etc.) \EE{dokumentiert} werden. Der
Geburtsort ist insofern wichtig, da für die Ausstellung der
Geburtsurkunde  das für den Geburtsort zuständige Bezirksamt
verantwortlich ist. Wenn Mutter und Kind wohlauf sind, kann
der Transport in ein Krankenhaus fortgesetzt werden.

Es soll nicht auf die Nachgeburt gewartet werden. Sollte die
Plazenta vor Eintreffen im Krankenhaus geboren werden, ist
diese unbedingt sicherzustellen und mitzunehmen. Sie wird im
Krankenhaus auf Vollständigkeit und Reife untersucht.



\subsection{Notfälle während und unmittelbar nach der Geburt}
\label{topic:notfall-geburt}

\subsubsection{Basisreanimation des Neugeborenen}

% M %%%%%%%%%%%%%%%%%%%%%%%%%%%%%%%%%%%%%%%%%%%%%%%%%%% M %
% Basisreanimation des Neugeborenen
\ProcedurePrintDefault{GY51140C}

\begin{PictureSeriesWide}{}{Basisreanimation des Neugeborenen}
\PictureSeriesRowWide{3}
{./images/ASBOE-Grassl-gjh-1/CRW_0611-00441pt}
{Beatmung mittels Beatmungsbeutel für Neugeborene \SourcePicture{ASBÖ/gjh}}
{./images/ASBOE-Grassl-gjh-1/CRW_0623-00441pt}
{Positionierung der Finger \SourcePicture{ASBÖ/gjh}}
{./images/ASBOE-Grassl-gjh-1/CRW_0627-00441pt}
{Reanimation \SourcePicture{ASBÖ/gjh}}
{}
{}
\end{PictureSeriesWide}

\subsubsection{Asphyxie während der Geburt}

\Text{\TxBeschreibung}{%

\DefinitionInPara{Eine \T{Asphyxie}
ist eine Sauerstoffunterversorgung (Hypoxie) des
Neugeborenen infolge einer atem- oder kreislaufbedingten 
Störung vor, während oder nach der Geburt.}
Die möglichen Ursachen für die Hypoxie sind vielfältig und
umfassen u.\,a. eine Placenta praevia, einen
Nabelschnurvorfall, Lage\-ano\-malien, eine Frühgeburt,
u.\,v.\,a.\,m. Wird die Ursache nicht rechtzeitig behoben
und bleibt die Hypoxie bestehen, kommt es zu einer
irreperablen Hirnschädigung.
}{%

}{}{}{}

\TextSlide{\TxSymptome}{%
Das Leitsymptom einer Asphyxie ist eine \EE{Herzfrequenz <
100\PerMin*}, eine \EE{unzureichende Atmung}
(Schnappatmung, fehlende Atmung) und/oder ein \EE{schlaffer
Muskeltonus} des Kindes
Weitere Hinweise auf  eine Asphyxie sind eine grünlich-trübe
Verfärbung des Fruchtwassers sowie eine \E{fortbestehende
Zyanose} (eine Zyanose unmittelbar nach der Geburt ist
normal).
}{%
\begin{itemize}
  \item HF < 100\PerMin*
  \item Unzureichende Atmung
  \item Schlaffer Muskeltonus
\end{itemize}
}{}{}{}


% M %%%%%%%%%%%%%%%%%%%%%%%%%%%%%%%%%%%%%%%%%%%%%%%%%%% M %
% Asphyxie des Neugeborenen
\ProcedurePrintDefault{GP21090C}

\subsubsection{Nabelschnurvorfall}

Bei einem Nabelschnurvorfall in den Geburtskanal bestehen
zwei Gefahren:
Einerseits wird während des Geburtsvorganges die Nabelschnur
durch den Kindskopf abgeklemmt, und somit ist die
Sauerstoffversorgung des Kindes nicht mehr gewährleistet.
Andererseits kann sich die Nabelschnur  um den Hals des
Kindes wickeln und es kommt während der Geburt zur
Strangulation des Kindes.

% M %%%%%%%%%%%%%%%%%%%%%%%%%%%%%%%%%%%%%%%%%%%%%%%%%%% M %
% Nabelschnurvorfall
\ProcedurePrintDefault{GO69000C}


\subsubsection{Placenta praevia}

\TextSlide{\TxBeschreibung}{%
Normalerweise liegt die Plazenta nicht vor dem Geburtskanal.
Bei einer sehr tiefen Einnistung des befruchteten Eis in der
Gebärmutter, kann es jedoch dazukommen, dass sich die
\EE{Plazenta \E{vor} dem Geburtskanal} bildet. Dies ist
natürlich ein sehr großes Geburtshindernis und i.\,d.\,R.
ist eine vaginale Geburt unmöglich. Außerdem kann es durch
vorzeitiges Lösen des Mutterkuchens zu massiven Blutungen
kommen.
Eine Placenta praevia ist meistens schon durch die
\E{Kontrolluntersuchungen} beim Frauenarzt bekannt und kann
erfragt bzw. im Mutter-Kind-Pass nachgelesen werden.
}{%
\begin{itemize}
  \item Plazenta liegt vor dem Geburtskanal
  \item Vaginale Geburt nicht möglich
  \item Blutungsgefahr 
  \item Anamnese, Mutter-Kind-Pass
\end{itemize}
}{}{}{}


\TextSlide{\TxSymptome}{%
Unter Wehentätigkeit kommt es zu starken Blutungen aus der
Scheide.

\vspace{0.5cm}

}{%
\begin{itemize}
  \item Vaginalblutungen während Wehentätigkeit
\end{itemize}
}{}{}{}

% M %%%%%%%%%%%%%%%%%%%%%%%%%%%%%%%%%%%%%%%%%%%%%%%%%%% M %
% Placenta praevia bei bevorstehender Geburt
\ProcedurePrintDefault{GO69041C}

\subsubsection{Pathologische Geburtslagen}


Alle Lagen des Kindes, bei denen nicht zuerst der 
Kopf im Geburtskanal erscheint:
\begin{itemize}
	\item Steißlage bzw. Beckenendlage
	\item Armvorfall
\end{itemize}
Die größte Gefahr bei allen Lageanomalien ist, dass die
Nabelschnur abgedrückt wird. Daher sollte unter allen
Umständen die Geburt im KH, mittels Kaiserschnitt,
durchgeführt werden. Die Geburt in jedem Fall verhindern
bzw. verzögern. Ein Armvorfall ist eine geburtsunmögliche
Lageanomalie. 

% M %%%%%%%%%%%%%%%%%%%%%%%%%%%%%%%%%%%%%%%%%%%%%%%%%%% M %
% Pathologische Geburtslagen
\ProcedurePrintDefault{GO64090C}


\subsubsection{Blutungen in der
Nachgeburtsperiode}

\TextSlide
{\TxBeschreibung}
{
% TODO Uterusatonie --> Blutungen in der Nachgeburtsperiode
Normalerweise zieht sich in der Nachgeburtsphase die
Uterusmuskulatur zusammen. In diesem Fall kommt es aber zur
Erschlaffung (Atonie) der Gebärmutter. Die Folge sind
lebensbedrohliche Blutungen durch das fehlende Kontrahieren
des Uterus nach der Plazentalösung.}
{
\begin{itemize}
  \item Erschlaffung der Gebärmutter
  \item Blutungen
\end{itemize}
}
{}
{}
{}


\MassnahmenMK{Spezielle
Maßnahmen}{Uterusatonie}{m:uterusatonie}{
}{}



\subsubsection{Uterusruptur}

\Text{\TxBeschreibung}{% 
Die Uterusruptur kommt meistens nur nach vorangegangenem
Kaiserschnitt vor. Es treten während der Wehen plötzlich
Schmerzen auf, welche anders als Wehenschmerzen empfunden
werden. Die Wehen nehmen dann ab oder hören plötzlich auf.
Damit wird der Geburtsvorgang gestoppt und es herrscht
Lebensgefahr für die Mutter aufgrund starker Blutungen und
für das Kind wegen des \ce{O2}-Mangels.
}{}{}{}{}

% M %%%%%%%%%%%%%%%%%%%%%%%%%%%%%%%%%%%%%%%%%%%%%%%%%%% M %
% Uterusruptur
\ProcedurePrintDefault{GO71010C}





% \clearpage % TEMP temp clearpage
\section{Sonstige gynäkologische Erkrankungen und
Notfälle}

\subsection{Vaginale Blutungen}
\label{topic:unterleibsblutungen}

\TextSlide{\TxBeschreibung}{%
Die vaginale Blutung ist präklinisch schwer abzuklären und
kann viele Ursachen haben. Wichtig ist festzuhalten, dass
jede vaginale Blutung abgeklärt gehört, da sich dahinter
auch ein Tumorgeschehen (besonders bei Patientinnen in der
Menopause) verstecken kann. Darüber hinaus kommt eine
verstärkte Monatsblutung, Verletzung (nach Sturz,
Geschlechtsverkehr oder nach Vergewaltigungen),
Defloration\footnote{Entjungferung} oder Blutungen im Rahmen
einer Schwangerschaft in Betracht. Bei vaginalen Blutungen
in Kindern darf man \E{nicht} automatisch von einem
erfolgten Missbrauch ausgehen.
}{
\begin{itemize}
    \item Verstärkte Monatsblutung
    \item Verletzungen (Sturz, Geschlechtsverkehr, Vergewaltigung,
    \ldots)
    \item Defloration
    \item Tumor
    \item Im Rahmen einer Schwangerschaft
\end{itemize}
}{}{}{}

Eine akute vitale Bedrohung besteht seltenst, dennoch muss
sie routinemäßig vor allem anhand des Blutverlustes und der
Vitalwerte eingeschätzt werden. 


% M %%%%%%%%%%%%%%%%%%%%%%%%%%%%%%%%%%%%%%%%%%%%%%%%%%% M %
% Vaginale Blutung
\ProcedurePrintDefault{GN93091C}


\begin{Literary}
  {\fullcite{HaagGyn2007}} 

  {\fullcite{BuehlingGynaekologie1}}

  {\fullcite{OHCSP7}}

  {\fullcite{Nilsson1997}} 
  
  Weiters: \cite{
SiewertChirurgie8,
AMLS3,
% Erc2005DeSec5,
% Erc2005Sec5,
% Erc2005DeSec6,
% Erc2005DeSec3,
% Erc2005Sec3,
BoehmerTaschRett6,
Fauci2008,
LPN3,
Boecker2,
% Helfen2008,
LippertAnatomie7,
RedelsteinerHRNS1,
ForMed2,
StriebelAn7,
KlinkePhysio3,
DRAn3,
ForthPharma8,
LuellmannPharma16,
AnCompact3,
RueckerBildatlas1}
\end{Literary}

% \section*{\protect\dsliterary{} Weiterführende Literatur}
% \begin{WidePar}
% \begin{multicols}{2}
% \Z{%
% \begin{labeling}{mmm}
%   \item[\cite{HaagGyn2007}]
%   \emph{\bibentry{HaagGyn2007}} 
%   \item[\cite{BuehlingGynaekologie1}] 
%   \emph{\bibentry{BuehlingGynaekologie1}}
%     \item[\cite{OHCSP7}] 
%   \emph{\bibentry{OHCSP7}}
%   \item[\cite{Nilsson1997}] 
%   \emph{\bibentry{Nilsson1997}} 
% \end{labeling}
% }
% \end{multicols}
% \end{WidePar}



\nocite{Juurlink2005,
Kalla2006,
Vlcek2008,
Wells2000,
Patzelt2005,
Sachsenweger2003,
}

\ChapterBibliography
\ChapterEnd
